\SourceLeaf{440}%
\[
\begin{aligned}
d\varphi={}&\frac{Z}{1-ZZ}(YdX-XdY)\\
&+\frac{d\xi}{1-ZZ}\{\zeta Y-\eta XXZ+\xi XYZ\}\\
&-\frac{d\eta}{1-ZZ}\{\zeta X-\eta XYZ+\xi YYZ\}\\
&+d\zeta(\eta X-\xi Y).
\end{aligned}
\]
Da nun ferner
\[
\begin{aligned}
\eta XXZ-\xi XYZ&=\eta XXZ+\eta YYZ+\zeta ZYZ\\
&=\eta Z(1-ZZ)+\zeta YZZ\\
\eta XYZ-\xi YYZ&=-\xi XXZ-\zeta XZZ-\xi YYZ\\
&=-\xi Z(1-ZZ)-\zeta XZZ,
\end{aligned}
\]
so wird unser ganzer Ausdruck
\[
\begin{aligned}
d\varphi={}&\frac{Z}{1-ZZ}(YdX-XdY)\\
&+(\zeta Y-\eta Z)d\xi+(\xi Z-\zeta X)d\eta+(\eta X-\xi Y)d\zeta.
\end{aligned}
\]

\Needspace{5\baselineskip}
\begin{center}15.\end{center}
\noindent\hspace*{1.2em}Die gefundene Formel gilt allgemein, wie auch die gekrümmte Linie beschaffen ist. Ist diese aber eine kürzeste Linie, so heben sich die drei letzten Terme offensichtlich auf, und folglich gilt
\[
d\varphi=-\frac{Z}{1-ZZ}(XdY-YdX).
\]
Man sieht aber leicht, dass
\[
\frac{Z}{1-ZZ}(XdY-YdX)
\]
nichts anderes ist als der Flächeninhalt des Teils der Hilfskugeloberfläche, der vom Element der Linie $L$, von den beiden Großkreisen durch seine Endpunkte und $(3)$ und vom dadurch abgeschnittenen Element des Großkreises durch $(1)$ und $(2)$ begrenzt wird. Dieser Flächeninhalt gilt als positiv, wenn $L$ mit $(3)$ auf derselben Seite von $(1)(2)$ liegt und die Richtung von $P$ nach $P'$ dieselbe ist wie von $(2)$ nach $(1)$, als negativ, wenn von einer dieser Bedingungen das Gegenteil gilt, und wieder als positiv, wenn von beiden das Gegenteil gilt. Mit anderen Worten ist er positiv, wenn man den
